\section{Background}
\label{sec:background}

\subsection{Cuckoo Filters}
\label{sec:background-cuckoo}

A Cuckoo Filter is an array of buckets, where each bucket contains several
slots that are read together \cite{fan2014cuckoo}.  Each occupied slot stores a
short fingerprint rather than a complete key.  Hashing a key produces its
fingerprint and two possible bucket locations.  A lookup checks both buckets,
while an insertion places the fingerprint in an empty slot or relocates an
existing fingerprint when both buckets are full.  The structure also supports
deletion.  A negative lookup proves absence, while a positive lookup remains a
possible match and must be checked against exact state.

For example, suppose the address of a cacheline produces fingerprint $f$ and
maps to buckets 3 and 11.  When the cacheline is filled into L2, the filter
stores $f$ in one of the four slots in either bucket.  A later L2 data request
recomputes the same fingerprint and bucket locations.  If neither bucket
contains $f$, the cacheline is not resident and the request can skip the L2 tag
and data array lookup before continuing along the miss path.  If either bucket
contains $f$, the request checks the exact L2 tag.  A tag hit reads the data
array, while a false positive simply follows the normal cache miss path.  The
fingerprint is deleted when the cacheline is evicted.  If an insertion fails,
the filter conservatively sends later requests to the exact L2 lookup.  The
filter therefore avoids unnecessary cache probes without changing correctness
\cite{moshovos2001jetty}.  It stores only metadata and has finite lookup and
update bandwidth.

\subsection{RDMA Request Processing}
\label{sec:background-rdma}

RDMA moves data between GPMs without involving a host processor in each
transfer \cite{kalia2016design,wei2018deconstructing,mittal2018revisiting}.
Each GPM's RDMA endpoint connects both the private L1 caches and the shared L2
to the wafer mesh.  The L1 facing path contains a request table and a pair of
TX and RX interfaces.  The request table retains every remote memory
transaction issued by local L1 caches until its response returns.  TX sends
these requests to remote GPMs, while RX returns incoming responses to the
corresponding L1 requests.  The L2 facing path provides another pair of TX and
RX interfaces.  RX forwards incoming remote requests to the local L2, and TX
sends the resulting data responses back to their requesting GPMs.  A bounded
request table limits the number of requester transactions in flight and
propagates backpressure to L1 when no entry is available.

Figure~\ref{fig:baseline-overview} follows a baseline remote request from the
requesting GPM's L1 to the owner GPM's L2.  At \ding{192}, address translation
identifies a remote page, and the L1 miss allocates a request table entry before
reaching TX.  At \ding{193}, requester TX injects the request packet into the
wafer mesh.  At \ding{194}, the owner L2 facing RX accepts the packet.  At
\ding{195}, RX forwards the recovered cacheline request to the owner L2.  An
owner L2 hit supplies the data directly, while an L2 miss accesses owner HBM.
The response returns through owner TX, the wafer mesh, and requester RX.  The
request table then completes the stored transaction and delivers the data to
the original L1 request.  In the baseline, same line misses from different
private L1 caches allocate separate entries and network packets before the
owner L2 can observe and merge them.

\subsection{HBM Organization and Row Locality}
\label{sec:background-hbm}

HBM provides high bandwidth through a wide interface and many independently
operated banks \cite{li2016exploring,wang2022shuhai,micron_hbm2e}, but each
bank still follows the conventional DRAM row buffer model
\cite{rixner2000memory,ding2013reshaping}.  A physical address is decomposed
into channel, rank, bank group, bank,
row, and column fields.  An activate command copies one row into the bank row
buffer.  A following column read to that open row can transfer data without
another activation.  A request to a different row must first precharge the
bank and then activate the new row.  Serving several nearby columns while a row
remains open can therefore avoid repeated precharge and activation latency
\cite{rixner2000memory}.  HBM address mapping and the order in which requests
reach the controller determine whether this locality is visible
\cite{li2016exploring,wang2022shuhai,wang2019mac}.  The simulator models these
effects explicitly with transaction and command queues, bank state, burst
timing, and separate activate, column, and precharge commands.  Its ordinary
close page path attaches an automatic precharge to a completed access.  When
two linked adjacent requests are simultaneously ready and map to the same
physical bank and row, the controller changes the first operation into a read
without automatic precharge, keeps the row open, and issues the adjacent column
next.  The row is not retained when the second request is unavailable, maps to
a different row, or would bypass an older conflicting request.  Adjacent line
selection is therefore useful because it creates a bounded opportunity for
row reuse, while the controller remains responsible for validating the actual
HBM mapping and timing constraints.
